\chapter{背景知识}

1. 什么是 Dify？

**Dify** 是一个**开源的大语言模型（LLM）应用开发平台**，旨在帮助开发者快速构建和部署基于 AI 的应用。它提供了一个直观的界面，集成了 AI 工作流、RAG（检索增强生成）、智能体（Agent）功能和模型管理等工具。简单来说，Dify 就像一个“搭积木”的平台，用户无需深入编码就能创建 AI 应用，例如智能客服或知识库问答系统。

 2. 什么是 Ollama？
 
**Ollama** 是一个**轻量级的本地推理框架**，用于在本地机器上部署和运行大语言模型（LLM）。它支持多种模型（如 LLaMA、Mistral、DeepSeek 等），用户只需通过简单命令即可下载和运行这些模型。Ollama 的优势在于本地化部署，不依赖云服务，保证数据隐私，同时提供 API 接口供其他应用调用。

 3. 什么是 Embedding？
 
**Embedding（嵌入）** 是一种**将文本、图像等数据转化为向量表示的技术**。在 AI 领域，Embedding 模型将复杂的自然语言或数据映射成数字向量，这些向量保留了语义信息。例如，“猫”和“狗”的向量会比较接近，而“猫”和“汽车”的向量距离较远。Embedding 在知识检索、语义搜索和问答系统中起着关键作用。

4. 什么是 LightRAG？

**LightRAG** 是一个**轻量级的 RAG（检索增强生成）系统**，结合了知识图谱和向量检索技术，以提高数据检索的效率和准确性。它通过将文本数据转化为知识图谱，并利用向量检索快速定位相关信息，为大语言模型提供更精准的上下文。LightRAG 注重轻量化和高效性。

5. 什么是 GraphRAG？

**GraphRAG** 是一种**基于图的检索增强生成方法**，利用知识图谱增强大语言模型的检索能力。它通过构建文本数据的知识图谱，并利用图中的社区结构和层次关系，生成更全面和准确的回答。GraphRAG 特别适合处理复杂查询和需要深入理解数据集上下文的场景。

 6. 什么是 LazyGraphRAG？
 
**LazyGraphRAG** 是一个**类似于 GraphRAG 的系统**，但在构建和查询知识图谱时采用了更高效的策略，以减少计算资源需求。它可能通过延迟加载或按需计算图谱的某些部分来优化性能，从而在保持准确性的同时提升效率。

7. 什么是 MinerU？

**MinerU** 是一个**知识图谱构建工具**，能够从文本数据中自动提取实体和关系，并构建知识图谱。它通常与 RAG 系统结合使用，提供结构化的知识表示，从而提高检索和生成的准确性。

- 8. 什么是 DeepSeek？

**DeepSeek** 是一个**开源的大语言模型**，在 RAG 系统中常用作生成模型。它以高性能和低成本著称，特别在处理大规模文档和复杂查询时表现出色。DeepSeek 可通过 Ollama 在本地运行，提供隐私保护和低延迟的 AI 服务。

 它们之间的联系
这些工具和技术在 AI 和大语言模型生态系统中相互关联，协同工作以构建高效、私密且功能强大的应用：

- **Dify** 作为一个开发平台，可以集成 **Ollama** 提供的本地模型服务。用户可在 Dify 中配置 Ollama 的 API 接口，使用本地运行的模型（如 **DeepSeek**），从而构建隐私性强的 AI 应用。
- **Embedding** 是 RAG 系统的核心技术，用于将文本数据转化为向量表示，支持 **LightRAG**、**GraphRAG** 和 **LazyGraphRAG** 的语义搜索和检索。Ollama 可提供本地 Embedding 模型（如 nomic-embed-text）。
- **LightRAG**、**GraphRAG** 和 **LazyGraphRAG** 是 RAG 系统的变体，目标是通过检索外部知识增强模型生成能力。**LightRAG** 强调高效性，**GraphRAG** 利用知识图谱提升准确性和深度，**LazyGraphRAG** 则优化资源消耗。
- **MinerU** 与 **GraphRAG** 和 **LazyGraphRAG** 结合，从文本中提取实体和关系，构建知识图谱，为检索提供结构化支持。
- **DeepSeek** 作为生成模型，可通过 Ollama 在本地运行，并集成到 Dify 应用中，结合 RAG 系统提供准确回答。

总结
- **Dify**：应用开发平台。
- **Ollama**：本地模型运行框架。
- **Embedding**：文本向量化技术。
- **LightRAG**：轻量级 RAG 系统。
- **GraphRAG**：基于图的 RAG 系统。
- **LazyGraphRAG**：高效优化的 GraphRAG。
- **MinerU**：知识图谱构建工具。
- **DeepSeek**：高性能生成模型。

这些工具通过协同工作，支持从知识提取到检索和生成的全流程，适用于本地化、私密性强的 AI 应用场景。



\section{1Panel}
\url{https://1panel.cn/}
1Panel 提供了一个直观的 Web 界面，帮助用户轻松管理 Linux 服务器中的网站、文件、容器、数据库以及大型语言模型（LLMs）。




